
\chapter{2016年上海卷（春）}

\section{填空题}

\begin{problem}
复数 \(3+4\i\)（\(\i\) 为虚数单位）的实部是\fillinblank{}.

\end{problem}

\begin{answer}
\(3\)
\end{answer}

\begin{problem}
若 \(\log_2(x+1)=3\)，则 \(x=\)\fillinblank{}.

\end{problem}

\begin{answer}
\(7\)
\end{answer}

\begin{problem}
直线 \(y=x-1\) 与直线 \(y=2\) 的夹角为\fillinblank{}.

\end{problem}

\begin{answer}
\(\dfrac{\pi}{4}\)
\end{answer}

\begin{problem}
函数 \(f(x)=\sqrt{x-2}\) 的定义域为\fillinblank{}.

\end{problem}

\begin{answer}
\([2,+\infty)\)
\end{answer}

\begin{problem}
三阶行列式 \(\begin{vmatrix} 1 & -3 & 5 \\ 4 & 0 & 0 \\ -1 & 2 & 1 \end{vmatrix}\)
中，元素 \(5\) 的代数余子式的值为\fillinblank{}.

\end{problem}

\begin{answer}
\(8\)
\end{answer}

\begin{problem}
函数 \(f(x)=\dfrac1x+a\) 的反函数的图像经过点 \((2,1)\)，则实数 \(a=\)\fillinblank{}.

\end{problem}

\begin{answer}
\(1\)
\end{answer}

\begin{problem}
在 \(\triangle ABC\) 中，若 \(A=30^\circ\)，\(B=45^\circ\)，\(BC=\sqrt6\)，则 \(AC=\)\fillinblank{}.

\end{problem}

\begin{answer}
\(2\sqrt3\)
\end{answer}

\begin{problem}
\(4\) 个人排成一排照相，不同排列方式的种数为\fillinblank{}（结果用数值表示）.

\end{problem}

\begin{answer}
\(24\)
\end{answer}

\begin{problem}
无穷等比数列 \(\{a_n\}\) 的首项为 \(2\)，公比为 \(\dfrac13\)，则 \(\{a_n\}\) 的各项的和为\fillinblank{}.

\end{problem}

\begin{answer}
\(3\)
\end{answer}

\begin{problem}
若 \(2+\i\)（\(\i\) 为虚数单位）是关于 \(x\) 的实系数一元二次方程 \(x^2+ax+5=0\) 的一个虚根，则 \(a=\)\fillinblank{}.

\end{problem}

\begin{answer}
\(-4\)
\end{answer}

\begin{problem}
函数 \(y=x^2-2x+1\) 在区间 \([0,m]\) 上的最小值为 \(0\)，最大值为 \(1\)，则实数 \(m\) 的取值范围是\fillinblank{}.

\end{problem}

\begin{answer}
\([1,2]\)
\end{answer}

\begin{problem}
在平面直角坐标系 \(xOy\) 中，点 \(A\)，\(B\) 是圆 \(x^2+y^2-6x+5=0\) 上的两个动点，且满足 \(\lvert AB\rvert=2\sqrt3\)，则 \(\lvert\overrightarrow{OA}+\overrightarrow{OB}\rvert\) 的最小值为\fillinblank{}.

\end{problem}

\begin{answer}
\(4\)
\end{answer}


\section{单选题}

\begin{problem}
满足 \(\sin\alpha>0\) 且 \(\tan\alpha<0\) 的角 \(\alpha\) 属于（\quad）.

\choices
    {第一象限}
    {第二象限}
    {第三象限}
    {第四象限}

\end{problem}

\begin{answer}
B
\end{answer}

\begin{problem}
半径为 \(1\) 的球的表面积为（\quad）.

\choices
    {\(\pi\)}
    {\(\dfrac43\pi\)}
    {\(2\pi\)}
    {\(4\pi\)}

\end{problem}

\begin{answer}
D
\end{answer}

\begin{problem}
在 \((1+x)^6\) 的二项展开式中，\(x^2\) 项的系数为（\quad）.

\choices
    {\(2\)}
    {\(6\)}
    {\(15\)}
    {\(20\)}

\end{problem}

\begin{answer}
C
\end{answer}

\begin{problem}
幂函数 \(y=x^{-2}\) 的大致图像是（\quad）.

\choices
    {\choicebitmap[width=0.15\paperwidth]{img/2016/shanghai_spring/q16_optA.png}}
    {\choicebitmap[width=0.15\paperwidth]{img/2016/shanghai_spring/q16_optB.png}}
    {\choicebitmap[width=0.15\paperwidth]{img/2016/shanghai_spring/q16_optC.png}}
    {\choicebitmap[width=0.15\paperwidth]{img/2016/shanghai_spring/q16_optD.png}}

\end{problem}

\begin{answer}
C
\end{answer}

\begin{problem}
已知向量 \(\bs a=(1,0)\)，\(\bs b=(1,2)\)，则向量 \(\bs b\) 在向量 \(\bs a\) 方向上的投影为（\quad）.

\choices
    {\(1\)}
    {\(2\)}
    {\((1,0)\)}
    {\((0,2)\)}

\end{problem}

\begin{answer}
A
\end{answer}

\begin{problem}
设直线 \(l\) 与平面 \(\alpha\) 平行，直线 \(m\) 在平面 \(\alpha\) 上，那么（\quad）.

\choices
    {直线 \(l\) 平行于直线 \(m\)}
    {直线 \(l\) 与直线 \(m\) 异面}
    {直线 \(l\) 与直线 \(m\) 没有公共点}
    {直线 \(l\) 与直线 \(m\) 不垂直}

\end{problem}

\begin{answer}
C
\end{answer}

\begin{problem}
用数学归纳法证明等式 \(1+2+3+\cdots+2n=2n^2+n\)（\(n\in\N^*\)）的第（ii）步中，假设 \(n=k\) 时原等式成立，那么在 \(n=k+1\) 时，需要证明的等式为（\quad）.

\choices
    {\(1+2+3+\cdots+2k+2(k+1)=2k^2+k+2(k+1)^2+(k+1)\)}
    {\(1+2+3+\cdots+2k+2(k+1)=2(k+1)^2+(k+1)\)}
    {\(1+2+3+\cdots+2k+(2k+1)+2(k+1)=2k^2+k+2(k+1)^2+(k+1)\)}
    {\(1+2+3+\cdots+2k+(2k+1)+2(k+1)=2(k+1)^2+(k+1)\)}

\end{problem}

\begin{answer}
D
\end{answer}

\begin{problem}
关于双曲线 \(\dfrac{x^2}{16}-\dfrac{y^2}{4}=1\) 与 \(\dfrac{y^2}{16}-\dfrac{x^2}{4}=1\) 的焦距和渐近线，下列说法正确的是（\quad）.

\choices
    {焦距相等，渐近线相同}
    {焦距相等，渐近线不相同}
    {焦距不相等，渐近线相同}
    {焦距不相等，渐近线不相同}

\end{problem}

\begin{answer}
B
\end{answer}

\begin{problem}
设函数 \(y=f(x)\) 的定义域为 \(\R\)，则“\(f(0)=0\)”是“\(y=f(x)\) 为奇函数”的（\quad）.

\choices
    {充分不必要条件}
    {必要不充分条件}
    {充要条件}
    {既不充分也不必要条件}

\end{problem}

\begin{answer}
B
\end{answer}

\begin{problem}
下列关于实数 \(a\)，\(b\) 的不等式中，不恒成立的是（\quad）.

\choices
    {\(a^2+b^2\ge2ab\)}
    {\(a^2+b^2\ge-2ab\)}
    {\(\left(\dfrac{a+b}{2}\right)^2\ge ab\)}
    {\(\left(\dfrac{a+b}{2}\right)^2\ge-ab\)}

\end{problem}

\begin{answer}
D
\end{answer}

\begin{problem}
设单位向量 \(\bs e_1\) 与 \(\bs e_2\) 既不平行也不垂直，对非零向量 \(\bs a=x_1\bs e_1+y_1\bs e_2\)，\(\bs b=x_2\bs e_1+y_2\bs e_2\) 有结论：

\begin{circlelist}
\item 若 \(x_1y_2-x_2y_1=0\)，则 \(\bs a\parallel\bs b\)；
\item 若 \(x_1x_2-y_1y_2=0\)，则 \(\bs a\perp\bs b\).
\end{circlelist}

关于以上两个结论，正确的判断是（\quad）.

\choices
    {\(\circled{1}\)成立，\(\circled{2}\)不成立}
    {\(\circled{1}\)不成立，\(\circled{2}\)成立}
    {\(\circled{1}\)成立，\(\circled{2}\)成立}
    {\(\circled{1}\)不成立，\(\circled{2}\)不成立}

\end{problem}

\begin{answer}
A
\end{answer}

\begin{problem}
对于椭圆 \(C_{(a,b)}:\dfrac{x^2}{a^2}+\dfrac{y^2}{b^2}=1\)（\(a>0\)，\(b>0\)，\(a\neq b\)），若点 \((x_0,y_0)\) 满足 \(\dfrac{x_0^2}{a^2}+\dfrac{y_0^2}{b^2}<1\)，则称该点在椭圆 \(C_{(a,b)}\) 内. 在平面直角坐标系中，若点 \(A\) 在过点 \((2,1)\) 的任意椭圆 \(C_{(a,b)}\) 内或椭圆 \(C_{(a,b)}\) 上，则满足条件的点 \(A\) 构成的图形为（\quad）.

\choices
    {三角形及其内部}
    {矩形及其内部}
    {圆及其内部}
    {椭圆及其内部}

\end{problem}

\begin{answer}
B
\end{answer}

\section{解答题}

\begin{problem}
如图，已知正三棱柱 \(ABC-A_1B_1C_1\) 的体积为 \(9\sqrt3\)，底面边长为 \(3\)，求异面直线 \(BC_1\) 与 \(AC\) 所成的角的大小.

\bitmapfigure[width=0.26\linewidth]{img/2016/shanghai_spring/q25_fig1.png}

\end{problem}

\begin{solution}
\(h=4\)，所以 \(\theta=\arccos\frac{3}{10}\).
\end{solution}

\begin{problem}
已知函数 \(f(x)=\sin x+\sqrt3\cos x\)，求 \(f(x)\) 的最小正周期及最大值，并指出 \(f(x)\) 取得最大值时 \(x\) 的值.

\end{problem}

\begin{solution}
最小正周期为 \(T=2\pi\). 当 \(x=\frac{\pi}{6}+2k\pi\)（\(k\in\Z\)）时，\(y_{\max}=2\).
\end{solution}

\begin{problem}
如图，汽车前灯反射镜与轴截面的交线是抛物线的一部分，灯口所在的圆面与反射镜的轴垂直，灯泡位于抛物线的焦点 \(F\) 处，已知灯口直径是 \(24\text{ cm}\)，灯深 \(10\text{ cm}\)，求灯泡与反射镜的顶点 \(O\) 的距离.

\bitmapfigure[width=0.30\linewidth]{img/2016/shanghai_spring/q27_fig1.png}

\end{problem}

\begin{solution}
由 \(y^2=14.4x\) 得 \(\lvert OF\rvert=3.6\text{ cm}\).
\end{solution}

\begin{problem}
已知数列 \(\{a_n\}\) 是公差为 \(2\) 的等差数列.

\begin{enumerate}
\item 若 \(a_1\)，\(a_3\)，\(a_4\) 成等比数列，求 \(a_1\) 的值；
\item 设 \(a_1=-19\)，数列 \(\{a_n\}\) 的前 \(n\) 项和为 \(S_n\). 数列 \(\{b_n\}\) 满足 \(b_1=1\)，\(b_{n+1}-b_n=\left(\dfrac12\right)^n\). 记 \(c_n=S_n+2^{n-1}b_n\)（\(n\in\N^*\)），求数列 \(\{c_n\}\) 的最小项 \(c_{n_0}\)（即 \(c_{n_0}\le c_n\) 对任意 \(n\in\N^*\) 成立）.
\end{enumerate}

\end{problem}

\begin{solution}
\begin{enumerate}
\item \(a_1=-8\).
\item \(c_n=n^2-20n+2^n-1\)；通过数列单调性定义，可得当 \(n=4\) 时，有 \(\{c_n\}\) 的最小项 \(c_4=-49\).
\end{enumerate}
\end{solution}

\begin{problem}
对于函数 \(f(x)\) 与 \(g(x)\)，记集合 \(D_{f>g}=\{x\mid f(x)>g(x)\}\).

\begin{enumerate}
\item 设 \(f(x)=2\lvert x\rvert\)，\(g(x)=x+3\)，求 \(D_{f>g}\)；
\item 设 \(f_1(x)=x-1\)，\(f_2(x)=\left(\dfrac13\right)^x+a\cdot3^x+1\)，\(h(x)=0\). 如果 \(D_{f_1>h}\cup D_{f_2>h}=\R\)，求实数 \(a\) 的取值范围.
\end{enumerate}

\end{problem}

\begin{solution}
\begin{enumerate}
\item \(D_{f>g}=(-\infty,-1)\cup(3,+\infty)\).
\item \(D_{f_1>h}=(1,+\infty)\). 由 \(D_{f_1>h}\cup D_{f_2>h}=\R\)，可得 \((-\infty,1]\subseteq D_{f_2>h}\)，即不等式 \(\left(\frac13\right)^x+a\cdot3^x+1>0\) 在 \((-\infty,1]\) 上恒成立. 令 \(t=\left(\frac13\right)^x\)，则 \(-a<t^2+t\) 在 \(\left[\frac13,+\infty\right)\) 上恒成立，即 \(-a<(t^2+t)_{\min}=\frac49\)，即 \(a>-\frac49\). 故所求范围为 \(\left(-\frac49,+\infty\right)\).
\end{enumerate}
\end{solution}

\section{附加题：选择题}

\begin{problem}
若函数 \(f(x)=\sin(x+\varphi)\) 是偶函数，则 \(\varphi\) 的一个值是（\quad）.

\choices
    {\(0\)}
    {\(\dfrac{\pi}{2}\)}
    {\(\pi\)}
    {\(2\pi\)}

\end{problem}

\begin{answer}
B
\end{answer}

\begin{problem}
在复平面上，满足 \(\lvert z-1\rvert=4\) 的复数 \(z\) 所对应的点的轨迹是（\quad）.

\choices
    {两个点}
    {一条线段}
    {两条直线}
    {一个圆}

\end{problem}

\begin{answer}
D
\end{answer}

\begin{problem}
已知函数 \(f(x)\) 的图像是折线段 \(ABCDE\)，如图，其中 \(A(1,2)\)，\(B(2,1)\)，\(C(3,2)\)，\(D(4,1)\)，\(E(5,2)\)，设 \(k,b\in\R\)，若直线 \(y=kx+b\) 与 \(f(x)\) 的图像恰有 \(4\) 个不同的公共点，则 \(k\) 的取值范围是（\quad）.

\bitmapfigure[width=0.36\linewidth]{img/2016/shanghai_spring/q36_fig1.png}

\choices
    {\((-1,0)\cup(0,1)\)}
    {\(\left(-\dfrac13,\dfrac13\right)\)}
    {\((0,1]\)}
    {\(\left[0,\dfrac13\right]\)}

\end{problem}

\begin{answer}
B
\end{answer}

\section{附加题：填空题}

\begin{problem}
椭圆的 \(\dfrac{x^2}{25}+\dfrac{y^2}{9}=1\) 的长半轴的长为\fillinblank{}.

\end{problem}

\begin{answer}
\(5\)
\end{answer}

\begin{problem}
已知圆锥的母线长为 \(10\)，母线与轴的夹角为 \(30^\circ\)，则该圆锥的侧面积为\fillinblank{}.

\end{problem}

\begin{answer}
\(50\pi\)
\end{answer}

\begin{problem}
小明用数列 \(\{a_n\}\) 记录某地区 2015 年 12 月份 \(31\) 天中每天是否下过雨，方法为：当第 \(k\) 天下过雨时，记 \(a_k=1\)；当第 \(k\) 天没下过雨时，记 \(a_k=-1\)（\(1\le k\le31\)）. 他用数列 \(\{b_n\}\) 记录该地区该月每天气象台预报是否有雨，方法为：当预报第 \(k\) 天有雨时，记 \(b_k=1\)；当预报第 \(k\) 天没有雨时，记 \(b_k=-1\)（\(1\le k\le31\)）. 记录完毕后，小明计算出 \(a_1b_1+a_2b_2+a_3b_3+\cdots+a_{31}b_{31}=25\)，那么该月气象台预报准确的总天数为\fillinblank{}.

\end{problem}

\begin{answer}
\(28\)
\end{answer}

\section{附加题：解答题}

\begin{problem}
对于数列 \(\{a_n\}\) 与 \(\{b_n\}\)，若对数列 \(\{c_n\}\) 的每一项 \(c_k\)，均有 \(c_k=a_k\) 或 \(c_k=b_k\)，则称数列 \(\{c_n\}\) 是 \(\{a_n\}\) 与 \(\{b_n\}\) 的一个“并数列”.

\begin{enumerate}
\item 设数列 \(\{a_n\}\) 与 \(\{b_n\}\) 的前三项分别为 \(a_1=1\)，\(a_2=3\)，\(a_3=5\)，\(b_1=1\)，\(b_2=2\)，\(b_3=3\)，若 \(\{c_n\}\) 是 \(\{a_n\}\) 与 \(\{b_n\}\) 的一个“并数列”，求所有可能的有序数组 \((c_1,c_2,c_3)\)；
\item 已知数列 \(\{a_n\}\)，\(\{c_n\}\) 均为等差数列，\(\{a_n\}\) 的公差为 \(1\)，首项为正整数 \(t\)；\(\{c_n\}\) 的前 \(10\) 项和为 \(-30\)，前 \(20\) 项和为 \(-260\)，若存在唯一的数列 \(\{b_n\}\)，使得 \(\{c_n\}\) 是 \(\{a_n\}\) 与 \(\{b_n\}\) 的一个“并数列”，求 \(t\) 的值所构成的集合.
\end{enumerate}

\end{problem}

\begin{solution}
\begin{enumerate}
\item \((1,3,5)\)、\((1,3,3)\)、\((1,2,5)\)、\((1,2,3)\).

\item 由题意，易得
\(\begin{cases}
a_n=t+n-1,\\
c_n=-2n+8
\end{cases}\)
（\(n,t\in\N^*\)）. 因为存在唯一的数列 \(\{b_n\}\)，使得 \(\{c_n\}\) 是 \(\{a_n\}\) 与 \(\{b_n\}\) 的一个“并数列”，所以
\(\begin{cases}
a_n\neq c_n,\\
b_n=c_n
\end{cases}\)
对一切正整数 \(n\) 恒成立，即 \(t+n-1\neq-2n+8\Rightarrow t\neq-3n+9\Rightarrow t\neq3\) 且 \(t\neq6\). 因此，满足条件的 \(t\) 的值所构成的集合为 \(\{t\mid t\neq3,\ t\neq6,\ t\in\N^*\}\).
\end{enumerate}
\end{solution}
